\section{Diagonal points on the new split crystalline component}
\label{sec:cubic-fixed-point-split}

The following argument extends AC-G55 to the remaining
split coefficient primes of AC-G54.  It uses the
actual fixed-type potentially crystalline property
only for a diagonal representation.  No assertion
that every point of the new component is ordinary
is needed.

The primary input for the coefficient-field issue
is Brinon--Conrad, \emph{CMI Summer School Notes on
$p$-adic Hodge Theory}, author version
\url{https://math.stanford.edu/~conrad/papers/notes.pdf},
Theorem~2.2.7, printed p.~15.
Its character clause assumes values in the
coefficient field of the local Galois group and
image of Lie dimension at most one.
We verify these hypotheses below.
Their Theorem~2.4.6, printed p.~20, and
Proposition~9.3.1, printed p.~146, supply respectively
the invariance of Hodge--Tate properties and
crystallinity under inertial restriction.
We do not extend the $\mathbb Q_p$-valued
classification in Proposition~8.3.4 to arbitrary
coefficient characters without proof.

\begin{lemma}[AC-G55-character: rank-one inertia with finite coefficients]
\label{AC-G55-character}
Let $L/\mathbb Q_\ell$ be finite and let
$a:G_{\mathbb Q_\ell}\to L^\times$ be a continuous
potentially crystalline character.
Use the convention that $\kappa_\ell$ has
cyclotomic exponent $1$, and let $h\in\mathbb Z$
be the unique exponent of the Hodge type of $a$.
There are an unramified character $u$ and a
finite-image character $t$, with values in
$\mathcal O_L^\times$, such that
\[
                  a=u\,t\,\kappa_\ell^h .
                                                               \tag{G55S.1}
\]
In particular
\[
             a|_{I_\ell}=t|_{I_\ell}\kappa_\ell^h,
       \qquad \operatorname{WD}(a)|_{I_\ell}=t|_{I_\ell}
                                                               \tag{G55S.2}
\]
with the covariant convention for the
Weil--Deligne inertial representation.
\end{lemma}
\begin{proof}
Compactness of the Galois group implies
$a(G_{\mathbb Q_\ell})\subset\mathcal O_L^\times$:
its valuation image is a compact subgroup of
the discrete torsion-free group $\mathbb Z$.
The filtered de Rham module of an $L$-linear
character of $G_{\mathbb Q_\ell}$ is a
one-dimensional $L$-space.
Its filtration has one jump, so, on restriction
of scalars to $\mathbb Q_\ell$, its Hodge exponent
is $h$ with multiplicity $[L:\mathbb Q_\ell]$.
Consequently
$b=a\kappa_\ell^{-h}$ is Hodge--Tate of weight zero
in every coefficient component.

Local reciprocity identifies the abelianized
Galois group with
$\widehat{\mathbb Z}\times\mathbb Z_\ell^\times$,
where the first factor is unramified.
Let $u$ be the unramified factor of $b$, and put
$t=b/u$.
Thus $t$ factors through $\mathbb Z_\ell^\times$;
its compact image is finite or has an open
subgroup isomorphic to $\mathbb Z_\ell$.
An unramified representation is Hodge--Tate
of weight zero: by the inertial restriction
statement in Brinon--Conrad Theorem~2.4.6 it
has the same weights as its trivial inertia
restriction.  Hence $t$ also has weight zero.

Choose a finite Galois extension
$F/\mathbb Q_\ell$ containing the normal closure
of $L$ in the chosen algebraic closure.
On restriction to $G_F$, the coefficient
idempotents in
$L\otimes_{\mathbb Q_\ell}\mathbb C_\ell$
are stable.
The identity-embedding summand is
$\mathbb C_\ell(t)$, and the weight-zero
Hodge--Tate comparison gives it a nonzero
$G_F$-invariant vector.
Here is the coefficient detail: the entire
restricted representation is a direct sum of
trivial semilinear $\mathbb C_\ell$-lines;
each coefficient idempotent commutes with $G_F$
and is therefore represented, in that comparison,
by a matrix over $F$.
Its rank-one image thus has a nonzero invariant.

The restricted character has values in
$\mathcal O_F^\times$, and its image still has
Lie dimension at most one.
The character clause of Brinon--Conrad
Theorem~2.2.7 therefore implies that
$t(I_F)$ is finite: an infinite inertia image
would make the indicated invariant space zero.
The subgroup $I_F$ has finite index in $I_\ell$,
so $t(I_\ell)$ is finite as well.
Under local reciprocity $I_\ell$ maps onto
$\mathbb Z_\ell^\times$, through which $t$
factors.  Hence the full image of $t$ is finite.
This proves (G55S.1) without any
$\mathbb Q_\ell$-valued restriction on $a$.

Both $u$ and $\kappa_\ell^h$ are crystalline
and have trivial Weil--Deligne inertia.
For the finite character $t$, take a finite
Galois extension killing it.
Its crystalline module there is
$F_0\otimes_{\mathbb Q_\ell}L(t)$.
Rational inertia acts trivially on the
unramified coefficient factor $F_0$ and by
$t$ on the representation factor.
This identifies its covariant inertial type
with $t|_{I_\ell}$ and proves (G55S.2).
\end{proof}

\begin{theorem}[AC-G55-split: the fixed-point field at every good coefficient prime]
\label{AC-G55-split}
Retain $e\geq1$, $\ell\ne2,3,269$ and every
source hypothesis in AC-G49--G54.
Use the actual prime of AC-G54 at a split
coefficient prime with nontrivial local $\chi$.
Then every actual local point $f$ satisfying
both \texttt{Upper} and \texttt{Lower} is the
cubic basepoint:
\[
             \operatorname{push}(f,r^{\rm univ})=r_\chi,
                        \qquad f=\overline x_0 .
                                                               \tag{G55S.3}
\]
Together with AC-G55 this proves
\texttt{PtHyp.fixed} for all good coefficient
primes in the constructed model.
Consequently AC-G55-M3 extends to that full
range, with its \texttt{M2} antecedent retained.
\end{theorem}
\begin{proof}
Only the new split case requires proof.
As in AC-G55, the two triangular conditions
make the representation diagonal on $H=G_K$.
The normalization at complex conjugation makes
its full semidirect matrix component diagonal.
The exact finite-coefficient and integrality
argument there places the three characters
$a_0,a_1,a_2$ in one finite extension $L$, with
the prescribed residual characters.

At the chosen place above $\ell$, the
representation is an actual point killing
$P_*$.
By property 3 of AC-G54 it is potentially
crystalline with the fixed types of its seed.
Since it is a direct sum of three characters,
each summand is potentially crystalline.
This follows also directly from the period
dimension criterion: each summand has period
dimension at most its rank, and the equality
for their direct sum forces equality for
each summand.
The Hodge exponents are the multiset
$\{e,0,-e\}$, and the Weil--Deligne inertia
characters are
\[
                   \{\psi_I,1,\psi_I^{-1}\},
                       \qquad \psi_I=\psi|_{I_\ell}.
                                                               \tag{G55S.4}
\]
The cubic factor is unramified at $\ell$ and
does not alter either multiset.

Apply AC-G55-character to the three summands.
For some assignment of exponents and finite
inertia characters this gives actual equalities
\[
 a_i|_{I_\ell}=\tau_i\kappa_\ell^{h_i},\qquad
 \{h_i\}=\{e,0,-e\},\qquad
 \{\tau_i\}=\{\psi_I,1,\psi_I^{-1}\}.
                                                               \tag{G55S.5}
\]
This step is where diagonalness is essential;
no actual-inertia equality was inferred
directly from a semisimplified type of a
nondiagonal representation.

The exact \texttt{GD.ine} and \texttt{GD.gen}
fields of
\path{Development/AutomorphicArithmetic/AutomorphicVarieties/HeightAndVanishing/Counts.lean}
give
\[
            \psi_I^2=1,\qquad
                         \ell-1\nmid2e .
                                                               \tag{G55S.6}
\]
Here the source writes \texttt{comp ell psi} for
the local Dirichlet factor.  Its identification
with $\psi_I$ follows from the literal
\texttt{Dat.Psi\_I} field in
\path{Development/AutomorphicArithmetic/AutomorphicVarieties/UnitaryThreeGroups/RhoKApply.lean},
which writes the inertia restriction of $\gamma$
as $\kappa_\ell^e$ times that factor.
Every $\tau_i$ therefore has order dividing two.
The middle residual character is $1$.
If $h_1=e$ or $-e$, reduction of (G55S.5)
and then squaring would give
$\overline\kappa_\ell^{\,2e}=1$ on inertia.
The mod-$\ell$ cyclotomic character is
surjective on $\mathbb F_\ell^\times$ there,
contradicting (G55S.6).
Thus $h_1=0$.
Now $\overline{\tau}_1=1$, and a root of
unity of order dividing two reduces injectively
when $\ell\ne2$, so $\tau_1=1$.

The remaining exponents are $e,-e$, and the
remaining inertia types are $\psi_I,\psi_I$.
The first residual character is
$\overline\psi_I\overline\kappa_\ell^{\,e}$.
Assigning it exponent $-e$ would again give
$\overline\kappa_\ell^{\,2e}=1$.
Hence the coordinates are forced:
\[
       a_0|_{I_\ell}=\gamma,\qquad
       a_1|_{I_\ell}=1,\qquad
       a_2|_{I_\ell}=\gamma^{-1}.
                                                               \tag{G55S.7}
\]
The other place of $K$ above $\ell$ is the
complex-conjugate one.
Both the actual representation and its seed
have matrix $I$ at that conjugation; its
semidirect action reverses and inverts the
diagonal characters.
Thus (G55S.7) holds at that place as well.

For clarity, the rest of the global argument
uses exactly the actual conditions of the
new ring.
Set $b=(\chi^j\gamma,1,\chi^j\gamma^{-1})$
on $H$ and $\delta_i=a_i/b_i$.
The $\delta_i$ reduce to $1$.
At any finite $w\nmid\ell$, let
$m=3\operatorname{ord}(\psi)$.
The retained \texttt{GD.ord} gives $\ell\nmid m$.
For $g\in I_w$, the power $g^m$ belongs
to the actual residual inertia kernel killed
by the ring.
Thus $a_i(g)^m=b_i(g)^m=1$.
A root of unity of order prime to $\ell$
that reduces to $1$ is $1$, so
$\delta_i(g)=1$.
At $w\mid\ell$, this equality follows from
(G55S.7).
Each $\delta_i$ is therefore unramified
at every finite place of $K$.
The norm-Euclidean argument in AC-G55 proves
class number one for $\mathbb Z[\omega]$;
global reciprocity then forces
$\delta_i=1$.
Normalization at conjugation proves equality
of the full representations.
Their universal entry values agree, and the
actual entry density and continuity proved
in AC-G51 and AC-G53 give $f=\overline x_0$.

Finally the proof of AC-G55-M3 uses only
this fixed-point field and the model's
\texttt{pts} field under \texttt{M2}.
The latter field is part of the all-good-prime
model in AC-G54-model.
The same torus-translate argument therefore
gives \texttt{M3} in the enlarged range,
with precisely the original antecedent.
\end{proof}

\paragraph{Scope and geometric dictionary.}
This closes the diagonal-point requirement for
the new split component.  It asserts no
ordinary flag for arbitrary nondiagonal points,
no mixed tangent vanishing and no numerical
Selmer bound.
The finite inertia character in (G55S.1) is
a finite local system, with the already
audited SGA~1, Expos\'e~V, \S\S4 and~7
finite-cover interpretation.
The Hodge and period conditions, and their
coefficient extension, are established by the
displayed local argument and the stated
primary $p$-adic input, not by that dictionary.
